\subsection*{The Real Axis}

\begin{definitionbox}{Definition}
\begin{definition}[Canonical Embedding of $\mathbb{R}$ into $\mathbb{C}$]
\label{def:canonical-embedding-r-into-c}
Define
\[
\iota:\mathbb{R}\to\mathbb{C},
\qquad
\iota(a):=[a,0]_{\mathbb{C}}.
\]
\end{definition}
\end{definitionbox}

\LeanFormalizes{def:canonical-embedding-r-into-c}{lra-lean}{LRA.VolumeII.ComplexNumbers.Construction.Construction}{of_real}{pending}

\begin{remark*}[Standard quantified statement]
\[
\iota:\mathbb{R}\to\mathbb{C}
\quad\text{and}\quad
\forall a\in\mathbb{R},\ \iota(a)=[a,0]_{\mathbb{C}}.
\]
\end{remark*}

\begin{remark*}[Interpretation]
The canonical embedding sends each real number to the complex class on the real axis with imaginary coordinate zero.
\end{remark*}

\begin{dependencies}
\begin{itemize}
  \item \hyperref[def:reals]{The Real Numbers}
  \item \hyperref[def:complex-numbers]{Complex Numbers}
  \item \hyperref[def:complex-class-notation]{Complex Class Notation}
\end{itemize}
\end{dependencies}

\begin{theorembox}{Theorem}
\begin{theorem}[The Real Embedding into $\mathbb{C}$ Is Injective]
\label{thm:embedding-r-into-c-injective}
If \(\iota(a)=\iota(b)\), then \(a=b\).
\hyperref[prf:embedding-r-into-c-injective]{\textit{Go to proof.}}
\end{theorem}
\end{theorembox}

\LeanFormalizes{thm:embedding-r-into-c-injective}{lra-lean}{LRA.VolumeII.ComplexNumbers.Construction.Construction}{of_real_is_injective}{pending}

\begin{remark*}[Standard quantified statement]
\[
\forall a,b\in\mathbb{R},\qquad
\iota(a)=\iota(b)\Longrightarrow a=b.
\]
\end{remark*}

\begin{remark*}[Predicate reading]
\[
\operatorname{Injective}(\iota)
\]
means that equal embedded complex classes come from equal real numbers.
\end{remark*}

\begin{remark*}[Interpretation]
The real-axis embedding loses no real-number information.
\end{remark*}

\begin{dependencies}
\begin{itemize}
  \item \hyperref[def:canonical-embedding-r-into-c]{Canonical Embedding of $\mathbb{R}$ into $\mathbb{C}$}
  \item \hyperref[def:complex-coordinate-equivalence]{Complex Coordinate Equivalence}
  \item \hyperref[thm:equality-transitivity]{Equality Transitivity}
\end{itemize}
\end{dependencies}

\begin{theorembox}{Theorem}
\begin{theorem}[The Real Embedding Preserves Addition and Multiplication]
\label{thm:embedding-r-into-c-preserves-operations}
For all \(a,b\in\mathbb{R}\),
\[
\iota(a+b)=\iota(a)+\iota(b),
\qquad
\iota(ab)=\iota(a)\iota(b).
\]
\hyperref[prf:embedding-r-into-c-preserves-operations]{\textit{Go to proof.}}
\end{theorem}
\end{theorembox}

\LeanFormalizes{thm:embedding-r-into-c-preserves-operations}{lra-lean}{LRA.VolumeII.ComplexNumbers.Construction.Construction}{of_real_preserves_multiplication}{pending}

\begin{remark*}[Standard quantified statement]
\[
\forall a,b\in\mathbb{R},\qquad
\iota(a+b)=\iota(a)+\iota(b)
\land
\iota(a\cdot b)=\iota(a)\cdot\iota(b).
\]
\end{remark*}

\begin{remark*}[Predicate reading]
\[
\operatorname{Homomorphism}(\iota,\mathbb{R},\mathbb{C},+)
\land
\operatorname{Homomorphism}(\iota,\mathbb{R},\mathbb{C},\cdot)
\]
means that $\iota$ preserves real addition and multiplication as complex addition and multiplication.
\end{remark*}

\begin{remark*}[Interpretation]
Arithmetic on the real axis agrees with the corresponding complex arithmetic.
\end{remark*}

\begin{dependencies}
\begin{itemize}
  \item \hyperref[def:canonical-embedding-r-into-c]{Canonical Embedding of $\mathbb{R}$ into $\mathbb{C}$}
  \item \hyperref[def:addition-on-c-classes]{Addition on $\mathbb{C}$-Classes}
  \item \hyperref[def:multiplication-on-c-classes]{Multiplication on $\mathbb{C}$-Classes}
  \item \hyperref[def:addition-on-r]{Addition on $\mathbb{R}$}
  \item \hyperref[def:multiplication-on-r]{Multiplication on $\mathbb{R}$}
\end{itemize}
\end{dependencies}

\begin{theorembox}{Theorem}
\begin{theorem}[The Real Embedding Preserves Zero and One]
\label{thm:embedding-r-into-c-preserves-zero-one}
\[
\iota(0)=0_{\mathbb{C}},
\qquad
\iota(1)=1_{\mathbb{C}}.
\]
\hyperref[prf:embedding-r-into-c-preserves-zero-one]{\textit{Go to proof.}}
\end{theorem}
\end{theorembox}

\LeanFormalizes{thm:embedding-r-into-c-preserves-zero-one}{lra-lean}{LRA.VolumeII.ComplexNumbers.Construction.Construction}{of_real_preserves_one}{pending}

\begin{remark*}[Standard quantified statement]
\[
\iota(0_{\mathbb{R}})=0_{\mathbb{C}}
\land
\iota(1_{\mathbb{R}})=1_{\mathbb{C}}.
\]
\end{remark*}

\begin{remark*}[Interpretation]
The real zero and one embed as the complex zero and one.


After these preservation results, it is standard to 
identify \(a\in\mathbb{R}\)
with \([a,0]_{\mathbb{C}}\in\mathbb{C}\). Under this convention,
\[
[a,b]_{\mathbb{C}}=[a,0]_{\mathbb{C}}+[0,b]_{\mathbb{C}}
=a+bi.
\]
\end{remark*}

\begin{dependencies}
\begin{itemize}
  \item \hyperref[def:canonical-embedding-r-into-c]{Canonical Embedding of $\mathbb{R}$ into $\mathbb{C}$}
  \item \hyperref[def:zero-one-i-in-c]{Zero, One, and the Imaginary Unit in $\mathbb{C}$}
  \item \hyperref[def:zero-in-r]{Zero in $\mathbb{R}$}
  \item \hyperref[def:one-in-r]{One in $\mathbb{R}$}
\end{itemize}
\end{dependencies}

\subsection*{Real and Imaginary Parts}

\begin{definitionbox}{Definition}
\begin{definition}[Real and Imaginary Parts]
\label{def:real-and-imaginary-parts}
For \(z=[a,b]_{\mathbb{C}}\), define
\[
\operatorname{Re}(z):=a,
\qquad
\operatorname{Im}(z):=b.
\]
\end{definition}
\end{definitionbox}

\begin{remark*}[Standard quantified statement]
\[
\forall z=[a,b]_{\mathbb{C}}\in\mathbb{C},\qquad
\operatorname{Re}(z)=a\land \operatorname{Im}(z)=b.
\]
\end{remark*}

\begin{remark*}[Interpretation]
The real and imaginary part maps recover the two real coordinates of a complex class.
\end{remark*}

\begin{dependencies}
\begin{itemize}
  \item \hyperref[def:complex-class-notation]{Complex Class Notation}
  \item \hyperref[def:reals]{The Real Numbers}
\end{itemize}
\end{dependencies}

\begin{theorembox}{Theorem}
\begin{theorem}[Real and Imaginary Parts Are Well-Defined]
\label{thm:real-imaginary-parts-well-defined}
The functions
\[
\operatorname{Re},\operatorname{Im}:\mathbb{C}\to\mathbb{R}
\]
are well-defined.
\hyperref[prf:real-imaginary-parts-well-defined]{\textit{Go to proof.}}
\end{theorem}
\end{theorembox}

\begin{remark*}[Standard quantified statement]
\[
\operatorname{WellDefined}(\operatorname{Re},\mathbb{C},\mathbb{R})
\land
\operatorname{WellDefined}(\operatorname{Im},\mathbb{C},\mathbb{R}).
\]
\end{remark*}

\begin{remark*}[Interpretation]
Equivalent representatives of a complex class have the same real and imaginary coordinates.


The construction of \(\mathbb{C}\) completes the standard ladder of number
systems in this volume:
\[
\mathbb{N}\subseteq\mathbb{W}\subseteq\mathbb{Z}\subseteq
\mathbb{Q}\subseteq\mathbb{R}\subseteq\mathbb{C}.
\]
The final step sacrifices order compatibility in exchange for algebraic power:
\(-1\) now has a square root.
\end{remark*}

\begin{dependencies}
\begin{itemize}
  \item \hyperref[def:complex-class-notation]{Complex Class Notation}
  \item \hyperref[def:reals]{The Real Numbers}
\end{itemize}
\end{dependencies}
